% !TEX program = pdflatex
% ==========================================================
% CLASE BEAMER MASTER+
% Fundamentos de Matemática Básica - Agronomía
% Clase 2: Notación Científica
% Universidad Santo Tomás - Talca
% ==========================================================

\newif\ifhandout
%\handoutfalse
\handouttrue

\ifhandout
\documentclass[aspectratio=169,10pt,handout]{beamer}
\else
\documentclass[aspectratio=169,10pt]{beamer}
\fi

% ----------------------------------------------------------
% PAQUETES
% ----------------------------------------------------------
\usepackage[spanish]{babel}
\usepackage[T1]{fontenc}
\usepackage[utf8]{inputenc}

\usepackage{lmodern}
\usepackage{amsmath,amsfonts,amssymb}
\usepackage{multicol}
\usepackage{booktabs}
\usepackage{array}
\usepackage{tikz}
\usepackage{fontawesome5}

% ----------------------------------------------------------
% COLORES
% ----------------------------------------------------------
\definecolor{USTBlue}{RGB}{0,70,127}
\definecolor{USTTeal}{RGB}{0,153,117}
\definecolor{USTGray}{RGB}{110,110,110}
\definecolor{USTLightBlue}{RGB}{230,240,248}
\definecolor{USTLightTeal}{RGB}{225,245,238}

% ----------------------------------------------------------
% TEMA
% ----------------------------------------------------------
\usetheme{Madrid}
\usecolortheme[named=USTBlue]{structure}
\setbeamertemplate{navigation symbols}{}

\setbeamercolor{frametitle}{bg=USTBlue,fg=white}
\setbeamercolor{block title}{bg=USTLightBlue,fg=USTBlue}
\setbeamercolor{block title example}{bg=USTLightTeal,fg=USTBlue}

% ----------------------------------------------------------
% COMANDOS
% ----------------------------------------------------------
\newcommand{\pp}{\pause}
\newcommand{\idea}[1]{\textcolor{USTBlue}{\textbf{#1}}}

% ----------------------------------------------------------
% DATOS
% ----------------------------------------------------------
\title{Fundamentos de Matemática Básica}
\subtitle{Unidad I: Operaciones algebraicas en los reales\\ Clase 4: Notación Científica}
\author{Prof. Luis González Alcaino}
\date{Universidad Santo Tomás — Agronomía}

% ==========================================================
\begin{document}
	
	% ----------------------------------------------------------
	\begin{frame}
		\titlepage
	\end{frame}
	
	% ----------------------------------------------------------
	\begin{frame}{Meta de la clase}
		
		\begin{block}{Objetivo}
			Utilizar la notación científica para representar números muy grandes o muy pequeños y aplicar operaciones básicas usando potencias de 10.
		\end{block}
		
		\pp
		
		\begin{block}{Lo que veremos}
			\begin{enumerate}
				\item ¿Por qué usar notación científica?
				\item Cómo escribir números en notación científica
				\item Cómo volver al número original
				\item Multiplicación y división
				\item Aplicaciones en Agronomía
			\end{enumerate}
			
		\end{block}
		
	\end{frame}
	
	% ----------------------------------------------------------
	\begin{frame}{¿Por qué usar notación científica?}
		
		\begin{block}{Problema}
			En ciencias aparecen números muy grandes o muy pequeños.
		\end{block}
		
		\pp
		
		Ejemplo:
		
		\[
		0.00000045
		\]
		
		\pp
		
		\[
		450000000
		\]
		
		\pp
		
		\begin{exampleblock}{Idea clave}
			La notación científica permite escribir estos números de forma clara y compacta.
		\end{exampleblock}
		
	\end{frame}
	
	% ----------------------------------------------------------
	\begin{frame}{Forma de la notación científica}
		
		\begin{block}{Definición}
			
			Un número en notación científica se escribe como:
			
			\[
			a \times 10^n
			\]
			
			donde:
			
			\begin{itemize}
				\item \(1 \le a < 10\)
				\item \(n\) es un número entero
			\end{itemize}
			
		\end{block}
		
		\pp
		
		\begin{exampleblock}{Ejemplo}
			
			\[
			3.5 \times 10^4
			\]
			
		\end{exampleblock}
		
	\end{frame}
	
	% ----------------------------------------------------------
	\begin{frame}{Ejemplo: número grande}
		
		Escribir:
		
		\[
		450000
		\]
		
		en notación científica.
		
		\pp
		
		Paso 1:
		
		\[
		4.5
		\]
		
		\pp
		
		Paso 2:
		
		Contar los lugares desplazados
		
		\[
		450000 = 4.5 \times 10^5
		\]
		
	\end{frame}
	
		\begin{frame}{Números mayores a uno escritos en notación científica}
				
				\begin{enumerate}
					
					\item $42\,200\,000 = 4,22 \times 10^{7}$
					
					\item $592\,000\,000\,000 = 5,29 \times 10^{11}$
					
					\item $2\,756\,300\,000 = 2,7563 \times 10^{9}$
					
					\item $244,885 = 2,44885 \times 10^{2}$
					
					\item La población del mundo se estima en $6\,800\,000\,000$ personas.  
					Escribir esta cantidad en notación científica.
					
					\[
					6,8 \times 10^{9}
					\]
					
				\end{enumerate}
				
				\vspace{0.3cm}
				
				\textbf{Observación:} 
				Si el número es mayor que $1$, la coma decimal se desplaza hacia la izquierda y el exponente $n$ es positivo.
							
		\end{frame}
		
	% ----------------------------------------------------------
	\begin{frame}{Ejemplo: número pequeño}
		
		Convertir
		
		\[
		0.00045
		\]
		
		\pp
		
		\[
		0.00045 = 4.5 \times 10^{-4}
		\]
		
		\pp
			\textbf{Observación:} 
		Si el número es \textbf{menor que 1}, la coma decimal se desplaza hacia la \textbf{derecha}. El exponente \(n\) será \textbf{negativo}.
		
		\begin{block}{Regla}
			
			\begin{itemize}
				
				\item mover decimal a la derecha → exponente negativo
				
				\item mover decimal a la izquierda → exponente positivo
				
			\end{itemize}
			
		\end{block}
		
	\end{frame}
	
	% ----------------------------------------------------------
	\begin{frame}{Volver al número original}
		
		\[
		3.2 \times 10^4
		\]
		
		\pp
		
		\[
		32000
		\]
		
		\pp
		
		\[
		5.6 \times 10^{-3}
		\]
		
		\pp
		
		\[
		0.0056
		\]
		
	\end{frame}
	
	% ----------------------------------------------------------
	\begin{frame}{Multiplicación}
		
		\[
		(3 \times 10^4)(2 \times 10^3)
		\]
		
		\pp
		
		Multiplicar coeficientes
		
		\[
		3 \times 2 = 6
		\]
		
		\pp
		
		Sumar exponentes
		
		\[
		10^4 \times 10^3 = 10^7
		\]
		
		\pp
		
		Resultado
		
		\[
		6 \times 10^7
		\]
		
	\end{frame}
	
	% ----------------------------------------------------------
	\begin{frame}{División}
		
		\[
		\frac{6\times10^8}{2\times10^3}
		\]
		
		\pp
		
		Dividir coeficientes
		
		\[
		6/2 = 3
		\]
		
		\pp
		
		Restar exponentes
		
		\[
		10^8 / 10^3 = 10^5
		\]
		
		\pp
		
		Resultado
		
		\[
		3\times10^5
		\]
		
	\end{frame}
	
	% ----------------------------------------------------------
	\begin{frame}{Aplicación Agronómica}
		
		\begin{block}{Ejemplo}
			
			Un fertilizante contiene:
			
			\[
			4\times10^{-3}
			\]
			
			gramos de micronutriente por gramo de producto.
			
		\end{block}
		
		\pp
		
		Si se aplican:
		
		\[
		10^3
		\]
		
		gramos de fertilizante.
		
		\pp
		
		\[
		(4\times10^{-3})(10^3)
		\]
		
		\pp
		
		Resultado:
		
		\[
		4
		\]
		
		gramos de micronutriente.
		
	\end{frame}
	
	% ----------------------------------------------------------
	\begin{frame}{Ejercicios}
		
		\begin{enumerate}
			
			\item Escribir en notación científica
			
			\[
			560000
			\]
			
			\item Escribir en notación científica
			
			\[
			0.00032
			\]
			
			\item Multiplicar
			
			\[
			(2\times10^3)(4\times10^2)
			\]
			
			\item Dividir
			
			\[
			\frac{8\times10^6}{2\times10^3}
			\]
			
		\end{enumerate}
		
	\end{frame}
	
	% ----------------------------------------------------------
	\begin{frame}{Respuestas}
		
		\begin{multicols}{2}
			
			\begin{enumerate}
				
				\item \(5.6\times10^5\)
				
				\item \(3.2\times10^{-4}\)
				
				\item \(8\times10^5\)
				
				\item \(4\times10^3\)
				
			\end{enumerate}
			
		\end{multicols}
		
	\end{frame}
	
	% ----------------------------------------------------------
	\begin{frame}{Ticket de salida}
		
		\begin{block}{Responde antes de terminar la clase}
			
			\begin{enumerate}
				
				\item Convertir
				
				\[
				0.000078
				\]
				
				\item Calcular
				
				\[
				(3\times10^2)(5\times10^3)
				\]
				
				\item ¿Por qué la notación científica es útil en ciencias?
				
			\end{enumerate}
			
		\end{block}
		
	\end{frame}
	
	% ----------------------------------------------------------
	\begin{frame}
		
		\centering
		
		{\Huge\color{USTTeal}\faSeedling}
		
		\vspace{0.3cm}
		
		{\Large Notación científica}
		
		\vspace{0.3cm}
		
		{\small Matemática para comprender fenómenos reales}
		
	\end{frame}
	
	\begin{frame}{Trabajo personal del estudiante}

	\centering

	\vspace{0.3cm}

	Libro Álgebra, trigonometría y geometría análitica
	Dennis G. Zill - Jaqueline.M Dewar
	
	\medskip
	
	Ejercicio 2.3 Páginas 68-70	

	\end{frame}	
	
\end{document}